% FinSim -- CDSL Reimagine Ideathon, Screening Round entry
% Compile with pdfLaTeX on Overleaf. Upload all .png files from this folder next to this file.
% A missing image is replaced by a labelled placeholder box, so the document compiles either way.
% Target: exactly 3 full pages. Tune with \figh below: lower it (e.g. 3.4cm) if it spills onto page 4,
% raise it (e.g. 4.3cm) if page 3 has empty space. Every image is a real screenshot of the FinSim prototype.

\documentclass[10pt,twocolumn]{article}

\usepackage[a4paper,margin=0.5in,columnsep=0.25in]{geometry}
\usepackage[T1]{fontenc}
\usepackage{newtxtext,newtxmath}
\usepackage{microtype}
\usepackage[dvipsnames]{xcolor}
\usepackage{graphicx}
\usepackage{booktabs}
\usepackage{tabularx}
\usepackage{array}
\usepackage{enumitem}
\usepackage{titlesec}
\usepackage[font=footnotesize,labelfont=bf,skip=2pt]{caption}
\usepackage[hidelinks]{hyperref}
\usepackage{qrcode} % QR code in the "Try the Prototype" section; delete that line and the \qrcode call if it ever fails to load

% ---- Colours taken from the prototype ------------------------------------------
\definecolor{brand}{HTML}{0E6B57}
\definecolor{coral}{HTML}{D9472F}
\definecolor{soft}{HTML}{F5F0E6}
\definecolor{mint}{HTML}{DCEFE7}

% ---- Compact, paper-like layout ---------------------------------------------------
\titleformat{\section}{\normalfont\large\bfseries\color{brand}}{\thesection}{0.5em}{}
\titleformat{\subsection}{\normalfont\normalsize\bfseries}{\thesubsection}{0.5em}{}
\titlespacing*{\section}{0pt}{6pt}{2pt}
\titlespacing*{\subsection}{0pt}{4pt}{1pt}
\setlist{nosep,leftmargin=1.1em,itemsep=1pt}
\setlength{\parskip}{2pt}
\setlength{\parindent}{0pt}
\setlength{\textfloatsep}{6pt}
\setlength{\floatsep}{5pt}
\setlength{\intextsep}{5pt}
\renewcommand{\arraystretch}{1.12}
\newcolumntype{L}{>{\raggedright\arraybackslash}X}

% Tables and figures are NOT floats: each is an unbreakable block that stays exactly where it is written.
\newenvironment{inlinetable}{\par\smallskip\noindent\begin{minipage}{\columnwidth}\centering\footnotesize}{\end{minipage}\par\smallskip}
\newenvironment{inlinefig}{\par\smallskip\noindent\begin{minipage}{\columnwidth}\centering}{\end{minipage}\par\smallskip}
\raggedbottom % never stretch a column to fill the page

% Height of the small in-column screenshot pairs. Lower it to save space.
\newcommand{\figh}{3.8cm}

% \shot{file}{width}{height}: image scaled to fit the box, or a placeholder if missing.
\newcommand{\shot}[3]{%
  \IfFileExists{#1}{\includegraphics[width=#2,height=#3,keepaspectratio]{#1}}{%
    \fbox{\parbox[c][#3][c]{\dimexpr#2-2\fboxsep-2\fboxrule\relax}{\centering\scriptsize\color{gray}\texttt{#1}}}}}

% \pair{file1}{caption1}{file2}{caption2}{figure caption}{label}: two small shots side by side in one column.
\newcommand{\pair}[6]{%
  \begin{inlinefig}
  \begin{minipage}[b]{0.49\columnwidth}\centering\shot{#1}{\linewidth}{\figh}\\[-1pt]{\scriptsize\itshape #2}\end{minipage}\hfill
  \begin{minipage}[b]{0.49\columnwidth}\centering\shot{#3}{\linewidth}{\figh}\\[-1pt]{\scriptsize\itshape #4}\end{minipage}
  \captionof{figure}{#5}\label{#6}
  \end{inlinefig}}

% Shaded pitch box.
\newcommand{\pitch}[1]{\par\noindent\colorbox{mint}{\parbox{\dimexpr\columnwidth-2\fboxsep}{\small #1}}\par\vspace{2pt}}

\begin{document}

% ================================================================================
\twocolumn[{%
\begin{center}
  {\LARGE\bfseries FinSim: Learn the Market, and the People Trying to Trick You}\\[3pt]
  {\normalsize Flintech \;$\cdot$\; Aditi Choudhary, Harshil Agarwal\;$\cdot$\; Mentor: Prof. Debashish Chakravarty\;$\cdot$\; IIT Kharagpur}\\[1pt]
  {\small CDSL Reimagine Ideathon 2027 \;$\cdot$\; Screening Round Submission}
\end{center}
\vspace{-3pt}
\colorbox{soft}{\parbox{\dimexpr\textwidth-2\fboxsep}{\small
\textbf{Abstract.} India is adding investors faster than it is adding \emph{informed} investors. Only $\sim$9.5\% of households
invest in securities, participation in SEBI's investor education programmes remains below 1\%, and 51\% of non-investors cite
trust as a key barrier, while fraudsters now reach people directly through WhatsApp, video calls and deepfakes.
\textbf{FinSim} answers with \emph{practice, not lectures}: a free, browser-based platform where people get scammed safely and
learn the market by living it. Six two-minute, shareable \textbf{scam games} (digital arrest, pension/KYC calls, deepfakes,
pump-and-dump groups, fake IPO allotments, and the ``golden hour'' after fraud) pull people in; an eight-chapter
\textbf{investor course} with virtual money and a virtual phone takes them from ``what is a share'' to reading their own CAS. The
games attract; the course educates. \textbf{The prototype already works, in English and Hindi, with read-aloud, large text
and a built-in before/after self-check.}}}
\vspace{6pt}
}]

% ================================================================================
\section{About the Idea}

\subsection{Idea title}
\textbf{FinSim: Learn the Market, and the People Trying to Trick You.}

\subsection{The problem: reach, not value}
\begin{itemize}
  \item \textbf{Low participation:} only 3.21\,cr of 33.72\,cr households ($\sim$9.5\%) invest in securities, though 63\% know at
        least one product \cite{ps}.
  \item \textbf{Education does not reach people:} participation in SEBI's investor education programmes remains below 1\%, yet
        91\% of past participants rate them highly useful \cite{ps}.
  \item \textbf{Trust is the barrier:} 51\% of non-investors cite trust as a key barrier; $\sim$80\% of households are risk-averse
        \cite{ps}.
  \item \textbf{The most exposed are the least prepared:} 18--29 year-olds score lower on financial literacy than older cohorts
        (NCFE), cyber fraud disproportionately affects seniors and new investors, and 80\% of demat accounts come from Tier II+
        cities \cite{ps}.
\end{itemize}
\pitch{\textbf{Our insight.} People do not lack information; they lack \emph{rehearsal}. Nobody reads a circular in the middle
of a ``police'' video call. FinSim lets people rehearse that moment, and dozens like it, before it happens for real.}

\subsection{Summary of the idea}
FinSim has two doors into the same goal. Everything inside is simulated: fictional companies, virtual money, no real accounts.

\textbf{(A) Scam games: the hook.} Six standalone, two-minute games (Table~\ref{tab:games}), each a link you can send on
WhatsApp. Players make choices in a realistic chat, call or feed, and every game ends with a debrief: a score, a review of each
choice, the red flags, and the real-world response (1930, cybercrime.gov.in, freezing accounts, SEBI SCORES) \cite{sebi}.

\pair{chat.png}{A fake ``inspector'' on a video call}{deepfake.png}{Real or deepfake? Judge, then learn why}
     {Scam games put the player inside the scam, with virtual consequences.}{fig:games}

\textbf{How a round plays.} Each game shows one situation at a time and asks for one decision. Choices cost or earn points, never
real money, and the story reacts: the fake officer pushes harder, a share price climbs, the minutes in The Golden Hour run
out. Within two minutes the player has felt the pressure, made a choice and seen exactly which tactic was used.

\begin{inlinetable}
\captionof{table}{The six scam games.}
\label{tab:games}
\begin{tabularx}{\columnwidth}{@{}p{1.95cm}Lp{1.5cm}@{}}
\toprule
\textbf{Game} & \textbf{What the player practises} & \textbf{Audience} \\
\midrule
The Golden Hour   & Acting fast in the first 60 min after fraud & All \\
Digital Arrest    & Leaving a fake ``police'' video call & Seniors, all \\
Real or Deepfake? & Judging cloned voices, fake ads, notices & All \\
The Pension Call  & Refusing OTP-based ``life certificate'' calls & Seniors \\
The Tip Group     & Resisting a pump-and-dump group & Students, investors \\
Red Flag Rush     & Spotting flags in IPO, app, KYC scams & Students, investors \\
\bottomrule
\end{tabularx}
\end{inlinetable}

\pair{pension.png}{Pension Call: the OTP really says ``add beneficiary''}{rush.png}{Red Flag Rush: tap the tricks against the clock}
     {Each game drills in financial safety habits}{fig:games2}

\textbf{(B) Investor course: the depth.} Eight chapters in three modules (Table~\ref{tab:course}). Every chapter
\emph{teaches through interaction first}, with explore cards, sorting, sequencing and hands-on labs (Fig.~\ref{fig:labs}), and
only then quizzes. Learners can start at any chapter.

\pair{pricelab.png}{Price lab: news moves buyers and sellers}{compound.png}{Compounding lab: time beats amount}
     {Concepts are learned by playing with them, not reading about them.}{fig:labs}

\begin{inlinetable}
\captionof{table}{Course structure: concept $\rightarrow$ interaction.}
\label{tab:course}
\begin{tabularx}{\columnwidth}{@{}p{2.45cm}L@{}}
\toprule
\textbf{Chapter} & \textbf{Concept $\rightarrow$ interaction} \\
\midrule
Money DNA              & Money personality $\rightarrow$ scenario quiz \\
What Is a Share?       & Owning vs lending $\rightarrow$ sort, price lab \\
Where Your Shares Live & Bank/trading/demat, CDSL/NSDL, DPs $\rightarrow$ sequence a trade \\
The Quiet Multiplier   & Compounding $\rightarrow$ slider lab, recurring deposit \\
Opening a Demat A/c    & SEBI registration, KYC, nominee, DDPI/TPIN $\rightarrow$ phone task \\
Reading Your CAS       & CAS anatomy $\rightarrow$ catch an unauthorised transfer \\
IPO, the Safe Way      & ASBA, UPI mandate, lottery $\rightarrow$ legit-or-scam sort \\
Investor Scam Radar    & Fake apps, tip groups, fake regulators $\rightarrow$ scam autopsy \\
\bottomrule
\end{tabularx}
\end{inlinetable}

\pair{anatomy.png}{Explore: tap each part of a CAS}{sort.png}{Sort: legit or scam, with instant feedback}
     {Teaching screens come before every quiz.}{fig:teach}

A \textbf{virtual phone} makes it real: users verify a broker in a simulated SEBI registration lookup, and later spot an
off-market transfer they never made in their own CAS (Fig.~\ref{fig:phone}).

\pair{verify.png}{Verify: ``StockGuru Pro'' is not registered}{cas.png}{CAS: an unauthorised transfer, flagged}
     {The virtual phone turns verification into a habit.}{fig:phone}

\subsection{How it empowers, includes and builds trust}
FinSim is built around the four values the problem statement asks for, and every feature maps to at least one of them
(Table~\ref{tab:values}).

\begin{inlinetable}
\captionof{table}{Alignment with the problem statement's four values.}
\label{tab:values}
\begin{tabularx}{\columnwidth}{@{}p{1.5cm}L@{}}
\toprule
\textbf{Value} & \textbf{How FinSim delivers it} \\
\midrule
Empowerment & Decide, see consequences on simulated money, learn why; independent decisions, never advice. \\
Trust       & Teaches \emph{verification}: SEBI registration, CAS as the source of truth, official channels only. \\
Safety      & Learn $\rightarrow$ recognise $\rightarrow$ avoid $\rightarrow$ respond, end to end. \\
Inclusion   & Scenarios per audience; English + Hindi, read-aloud, large text; browser-only, no sign-up. \\
\bottomrule
\end{tabularx}
\end{inlinetable}

\subsection{Has this been done before?}
Only in parts. \textbf{FiAlly} (India) is a financial life simulator for young people, focused on banking and life skills.
\textbf{SEBI's} programmes and Saarthi app and \textbf{RBI's} campaigns are largely informational. \textbf{Bad News} and
\textbf{Go Viral!} (University of Cambridge and partners, UK) proved that playing through a manipulation builds resistance to it
\cite{roozenbeek,basol}, but not for finance. \textbf{FinSim is, to our knowledge, the first to combine securities-market
literacy, investor-specific fraud practice and a response playbook for students, first-time investors and seniors together.}

\subsection{Unique aspects}
\begin{itemize}
  \item \textbf{Securities-native:} demat, depositories, DPs, CAS, DDPI/TPIN, ASBA and SEBI checks, not generic bank tips.
  \item \textbf{Safe failure:} fall for a scam, lose virtual savings, then see exactly how it was done.
  \item \textbf{Two doors:} two-minute games for reach, a deep course for literacy.
  \item \textbf{Viral by design:} every game is a challenge link (``I scored 80, can you?'') carrying no personal data.
  \item \textbf{Teach, then test;} start anywhere.
\end{itemize}

\begin{inlinefig}
\shot{home.png}{\columnwidth}{4.4cm}
\captionof{figure}{Working prototype: the course and all six games on one screen.}
\label{fig:home}
\end{inlinefig}

\subsection{Implementation}
\textbf{Technology.} A lightweight React/TypeScript web app; content is stored as data, so new scams and new languages need no
new code. Static hosting keeps running costs close to zero.

\textbf{Data protection and privacy by design.}
\begin{itemize}
  \item \textbf{Nothing to steal:} no login, no names, no phone numbers; it never asks for PAN, Aadhaar, OTPs or account details.
  \item \textbf{Stays on the device:} progress and scores live only in the user's browser, and one button wipes them.
  \item \textbf{Clean sharing:} share links carry only the game's name, never a score history or identity.
  \item \textbf{Measurement without tracking:} the validation phase uses anonymous, aggregate counts and opt-in surveys, consistent with
        data minimisation under the DPDP Act, 2023 \cite{dpdp}.
\end{itemize}

\textbf{Accessibility and inclusion, already built.}
\begin{itemize}
  \item \textbf{English and Hindi} with one tap: the home page, all six games and the self-check (Fig.~\ref{fig:inclusion}).
  \item \textbf{Read-aloud} for low-literacy users and seniors: every chat line, card and result is spoken in Hindi or English.
  \item \textbf{Large-text mode} that scales the whole site, and it works on any phone browser with no install.
  \item \textbf{Built-in before/after self-check}: two equivalent six-message tests measure learning on the device itself.
\end{itemize}

\pair{hindi.png}{The Pension Call, fully in Hindi}{selfcheck.png}{Before/after self-check (example result)}
     {Inclusion and measurement built into the prototype}{fig:inclusion}

\begin{inlinetable}
\captionof{table}{Roll-out plan.}
\label{tab:plan}
\begin{tabularx}{\columnwidth}{@{}p{1.45cm}L@{}}
\toprule
\textbf{Phase} & \textbf{Deliverables} \\
\midrule
1. Done          & Working prototype: 6 games, 8 chapters, 10 articles; English + Hindi, read-aloud, large text, before/after self-check. \\
2. Validate  & Real users at 2--3 colleges and a senior citizens' association; impact measured by the built-in self-check; Hindi course. \\
3. Scale         & Links in DP onboarding and CAS emails, SEBI programmes, colleges, WhatsApp; more regional languages. \\
\bottomrule
\end{tabularx}
\end{inlinetable}

\pitch{\textbf{Why CDSL is the perfect channel.} Every investor already receives a CAS. One line in that email (``Can you spot
the scam? Play in 2 minutes'') could put FinSim in front of millions of demat holders at almost no cost.}

% ================================================================================
\section{Challenges, Risks \& Assumptions}

\subsection{Potential roadblocks}
\begin{inlinetable}
\captionof{table}{Risks and mitigations.}
\label{tab:risks}
\begin{tabularx}{\columnwidth}{@{}p{2.1cm}L@{}}
\toprule
\textbf{Risk} & \textbf{Mitigation} \\
\midrule
Low digital literacy & Short games, large text, voice, play with family \\
New scams appear     & Monthly content updates from official advisories \\
Accuracy, compliance & Expert review; fictional entities; no advice or trading \\
Drop-off             & Two-minute games, no sign-up, shareable scores \\
Proving impact       & Built-in before/after self-check; anonymous aggregate metrics \\
\bottomrule
\end{tabularx}
\end{inlinetable}

\textbf{Assumptions:} users have basic smartphone and internet access; partners (colleges, DPs, associations) will share links;
domain experts can review content before scale-up.

\pair{tipgroup.png}{The Tip Group: the rocket before the dump}{golden.png}{Golden Hour: every move costs minutes}
     {Games cover the whole arc, from temptation to response.}{fig:more}

\subsection{Why will this idea work?}
\begin{itemize}
  \item \textbf{It is evidence-based:} inoculation games measurably improve people's ability to recognise manipulation
        \cite{roozenbeek,basol}.
  \item \textbf{It fixes the real bottleneck:} the gap is reach, not value. A two-minute link removes the seminar barrier.
  \item \textbf{It spreads itself:} a score you can challenge your family with is a reason to share; every share is a new learner.
  \item \textbf{It already works:} the prototype is built and tested, in two languages, with measurement included.
\end{itemize}

\begin{inlinetable}
\captionof{table}{FinSim against the evaluation criteria.}
\label{tab:criteria}
\begin{tabularx}{\columnwidth}{@{}p{1.75cm}L@{}}
\toprule
\textbf{Criterion} & \textbf{Evidence in FinSim} \\
\midrule
Innovation       & Scam inoculation games + securities simulation in one platform \\
UX / UI          & One-screen home; two-minute games; teach-then-test chapters \\
Behavioural impact & Rehearses real decisions; measured by the before/after self-check \\
Feasibility      & Working prototype today; static hosting; data-driven content \\
Scalability      & Shareable links; CAS-email and DP distribution; near-zero cost \\
Accessibility    & English + Hindi, read-aloud, large text; browser-only, no sign-up \\
Data protection  & No personal data collected; all progress stays on device \\
Presentation     & Working demo of every feature shown in this paper \\
\bottomrule
\end{tabularx}
\end{inlinetable}

% ================================================================================
\section{Impact}

\subsection{Who will benefit the most}
\begin{itemize}
  \item \textbf{Senior citizens:} recognise digital-arrest, pension and impersonation scams; know whom to call.
  \item \textbf{First-time investors}, especially in Tier II+ cities: practical command of demat, CAS, IPOs and investor fraud.
  \item \textbf{Students (18--29):} resistance to finfluencer tips, fake trading apps and pump-and-dump groups.
\end{itemize}

\pitch{\textbf{Three learners, three journeys} (illustrative).\\[1pt]
\textbf{Kamala, 68, retired teacher:} her grandson forwards \emph{The Pension Call}. She learns to read what an OTP is for. When a
real ``treasury officer'' calls, she hangs up and phones her branch.\\[1pt]
\textbf{Rohan, 20, student:} added to a stock-tip group, he plays \emph{The Tip Group} first, checks the ``SEBI-registered expert''
and exits. Then he challenges his friends.\\[1pt]
\textbf{Priya, 29, first-time investor in Indore:} takes the course before opening her demat account, adds a nominee, never
shares her TPIN, and now reads every CAS.}

\begin{inlinefig}
\shot{debrief.png}{\columnwidth}{2.6cm}
\captionof{figure}{Every game ends with a score and a challenge to share.}
\label{fig:debrief}
\end{inlinefig}

\subsection{Expected outcomes and benefits}
\begin{inlinetable}
\captionof{table}{Outcomes and how they will be measured.}
\label{tab:kpi}
\begin{tabularx}{\columnwidth}{@{}LL@{}}
\toprule
\textbf{Outcome} & \textbf{Measure} \\
\midrule
Better scam detection    & Built-in before/after self-check score \\
Knows the response steps & \% who can name 1930 and account-freeze steps \\
Verification habits      & \% who check SEBI registration / read CAS \\
Reach                    & Plays, completions, shares (anonymous) \\
Confidence in markets    & Self-reported confidence, before/after \\
\bottomrule
\end{tabularx}
\end{inlinetable}

The result is the investor the problem statement calls for: informed, responsible, confident, and \emph{Atmanirbhar}.

\pitch{\textbf{Our goal:} help people make safer financial decisions \emph{before} a scammer gets the chance to make them for
real. }

% ================================================================================
\section{Try the Prototype}

The working prototype is live. It needs no sign-up and no install, and it runs in any phone or laptop browser.

\par\smallskip\noindent
\begin{minipage}[c]{0.63\columnwidth}\raggedright
\textbf{Website:}\\[1pt]
\href{https://dukenukem1611.github.io/cdsl-ideathon/}{\textcolor{brand}{\texttt{dukenukem1611.github.io/cdsl-ideathon}}}\\[4pt]
\textbf{A 60-second tour:} (1) tap \emph{The Digital Arrest} and play it; (2) switch the site to Hindi with the button at the top;
(3) read the debrief at the end; (4) open \emph{Self-check} and take the before test.
\end{minipage}\hfill
\begin{minipage}[c]{0.33\columnwidth}\centering
\qrcode[height=2.3cm]{https://dukenukem1611.github.io/cdsl-ideathon/}\\[1pt]{\scriptsize Scan to open}
\end{minipage}
\par\smallskip

% ================================================================================
\begingroup
\footnotesize
\begin{thebibliography}{9}\setlength{\itemsep}{0pt}
\bibitem{ps} CDSL, \emph{Reimagine Ideathon: Problem Statement}, 2026. \url{https://ideathon.cdslindia.com}
\bibitem{roozenbeek} J.~Roozenbeek and S.~van der Linden, ``Fake news game confers psychological resistance against online
  misinformation,'' \emph{Palgrave Communications}, 5, 65, 2019.
\bibitem{basol} M.~Basol, J.~Roozenbeek and S.~van der Linden, ``Good news about Bad News: Gamified inoculation boosts confidence
  and cognitive immunity against fake news,'' \emph{Journal of Cognition}, 3(1), 2, 2020.
\bibitem{dpdp} Government of India, \emph{Digital Personal Data Protection Act}, 2023.
\bibitem{sebi} SEBI, \url{https://www.sebi.gov.in}; SCORES, \url{https://scores.sebi.gov.in}; National Cyber Crime Reporting
  Portal, \url{https://cybercrime.gov.in} (helpline 1930).
\end{thebibliography}
\endgroup

\end{document}
